\documentclass[10pt,a4paper]{amsart}
\usepackage[margin=19mm]{geometry}
\usepackage{amsmath,amssymb,mathtools}
\usepackage[scr=rsfs,cal=euler]{mathalfa}
\usepackage{xfrac}
\usepackage[unicode,hidelinks,bookmarks=false]{hyperref}
\newcommand{\ie}{i.e.\ }
\newcommand{\cf}{cf.\ }
\newcommand{\resp}{respectively\ }
\newcommand{\Subsection}{\subsection}
\providecommand{\nobreakdash}{\nobreak\hbox{-}}
\setlength{\emergencystretch}{2em}
\multlinegap0pt
\usepackage{amssymb}
\usepackage{dsfont}
\newtheorem{theorem}{Theorem}
\newcommand {\bCP} {\mathbb {CP}}
\title{Plane rectifiable curves: old and new}
\author{Boris Shapiro, Guillaume Tahar}
\date{Source: arXiv:2605.09626v1 (CC BY 4.0)}
\begin{document}
\maketitle

\noindent\emph{Source and licence.} Plane rectifiable curves: old and new. Version arXiv:2605.09626v1 (CC BY 4.0), distributed by arXiv under the Creative Commons Attribution 4.0 International licence, http://creativecommons.org/licenses/by/4.0/, as recorded in the arXiv licence metadata for this exact version and shown on the abstract page https://arxiv.org/abs/2605.09626v1. Full licence notice: \texttt{licenses/arxiv-2605.09626-CC-BY-4.0.txt}.\par\smallskip

\noindent\textit{Editorial scope.} Counted unit: the article's theorem thm:two-singularities-circle (source0.tex lines 744--754). The article's complete original proof of it is delivered with it (source0.tex lines 756--795, one \texttt{proof} environment of 40 lines). No mathematics is written, repaired or re-derived: the statement and the proof are the article's own lines, in the article's own notation, and no retained line is substituted or re-flowed.  No further source-local context is needed: every cross-reference of every retained line resolves inside the two delivered spans.  Nothing was omitted from any retained span, so the package carries no omission record.  No retained line contains a citation, so the wrapper carries no bibliography and no bibliography entry.  No retained line contains a figure environment, an image-inclusion command or drawing code, so no image is delivered and the package declares no asset dependency.  Editorial formatting only, in addition: an AMS \texttt{amsart} preamble that restates the article's own declarations and macros used by the retained lines, namely the environments \texttt{theorem} on separate counters, together with the standard packages those lines need.  The retained environments print in this document as Theorem 1 in order of appearance, whereas the article prints the same items with the numbering of its own full document.  The complete native source of the article as distributed in the arXiv e-print for this exact version is bundled unchanged as \texttt{sources/200242/source0.tex}.
\begin{theorem}\label{thm:two-singularities-circle}
Let $\gamma:\bCP^{1}\to \bCP^{2}$ be a non-constant rational real plane curve, and let
\[
q=dx^{2}+dy^{2}
\]
be its Euclidean arc length quadratic differential on the normalization. Assume that $q$ has exactly two singularities. If the image of $\gamma$ is not a line, then it is a circle. More precisely, after a real M\"obius change of parameter and a Euclidean similarity, the complex coordinate $z=x+iy$ has the form
\[
z(t)=z_{0}+c\left(\frac{t+i}{t-i}\right)^{m},
\]
where $c\neq 0$ and $m\neq 0$ is an integer. In particular, if the parametrization is birational onto its image, then $|m|=1$.
\end{theorem}

\begin{proof}
Put
\[
z=x+iy,\qquad w=x-iy .
\]
Then $q=dz\,dw$, and $w(t)=\overline{z(\bar t)}$. In particular, the divisor of $q$ is invariant under the real structure on the parameter sphere.

We first discuss the case in which the two singularities are real. After a real M\"obius change of parameter, they are $0$ and $\infty$. Since, at a real point, $dz$ and $dw$ have the same order, the orders of $q$ at both singularities are even. Thus
\[
q=Ct^{2a}dt^{2}
\]
for some integer $a$ and some $C\neq0$. Correspondingly, $dz=\lambda t^{a}dt$ after changing the non-zero constant $\lambda$. If $a=-1$, then $z$ has a logarithmic term, impossible for a rational parametrization. If $a\neq -1$, then
\[
z=z_{0}+ct^{a+1}.
\]
Since $w=\overline{z(\bar t)}$, the real parametrization is contained in the real affine line determined by the complex direction $c$. Hence in the real case the image is a line.

It remains to consider the case in which the two singularities are non-real conjugate points. Since the divisor of a quadratic differential on $\bCP^1$ has degree $-4$, the two conjugate singularities have the same order and therefore both are double poles. After a real M\"obius change of parameter, they are $i$ and $-i$, and
\[
q=C\frac{dt^{2}}{(t^{2}+1)^{2}}
\]
for some $C\neq0$. Since $q=dz\,dw$, there is an integer $a$ and non-zero constants $\lambda,\mu$ such that
\[
dz=\lambda\frac{(t+i)^{a}}{(t-i)^{a+2}}dt,\qquad
 dw=\mu\frac{(t-i)^{a}}{(t+i)^{a+2}}dt .
\]
Set
\[
s=\frac{t+i}{t-i}.
\]
Then $s$ maps the real line to the unit circle and
\[
dz=-\frac{\lambda}{2i}s^{a}ds .
\]
The case $a=-1$ would give a logarithm, so it is impossible for a rational parametrization. Therefore
\[
z=z_{0}-\frac{\lambda}{2i(a+1)}s^{a+1}.
\]
For real $t$ one has $|s|=1$, so the image is contained in a Euclidean circle. Since the curve is irreducible and non-constant, its image is this circle. Taking $m=a+1$ gives the asserted normal form. Finally, a birational parametrization of the circle has degree one, so $|m|=1$.
\end{proof}

\end{document}
